% BHU PYQ -- Professional Exam Template
\documentclass[11pt,a4paper]{article}

\usepackage[a4paper,top=2.5cm,bottom=2.5cm,left=2.8cm,right=2.8cm,headheight=14pt]{geometry}
\usepackage{amsmath,amssymb}
\usepackage[T1]{fontenc}
\usepackage[utf8]{inputenc}
\usepackage{lmodern}
\usepackage{microtype}
\usepackage{fancyhdr}
\usepackage{enumitem}
\usepackage{hyperref}
\usepackage{lastpage}
\usepackage{parskip}

\hypersetup{colorlinks=true,urlcolor=black,linkcolor=black,
  pdftitle={B.Sc. (Hons.) Zoology Sem-V ZOB-505 -- BHU PYQ}}

\pagestyle{fancy}
\fancyhf{}
\renewcommand{\headrulewidth}{0.4pt}
\renewcommand{\footrulewidth}{0.4pt}
\fancyhead[L]{\small   \textit{Banaras Hindu University}}
\fancyhead[R]{\small   \textit{Zoology Paper: ZOB-505}}
\fancyfoot[C]{\small Page  \thepage\ of \pageref{LastPage}}


\newlist{parts}{enumerate}{3}
\setlist[parts,1]{label=(\alph*),leftmargin=*,itemsep=4pt,topsep=3pt}
\setlist[parts,2]{label=(\roman*),leftmargin=*,itemsep=2pt,topsep=2pt}
\setlist[parts,3]{label=(\arabic*),leftmargin=*,itemsep=2pt,topsep=2pt}

\newcommand{\pts}[1]{\hfill   \textit{[#1]}}


\begin{document}


\begin{center}
  {\large\bfseries BANARAS HINDU UNIVERSITY}\\[2pt]
  {\normalsize B.Sc. (Hons.) Semester V Examination, 2016--17}\\[6pt]
  \rule{0.8\linewidth}{0.5pt}\\[6pt]
  {\large\bfseries Zoology}\\[4pt]
  {\large\bfseries Paper No. ZOB-505 : Environmental Biology \& Systematics}\\[4pt]
  {\normalsize Time: Three hours \hfill Max. Marks: 70}
\end{center}

\rule{\linewidth}{0.4pt}

\smallskip



\noindent   \textit{Note: Write your Roll No. at the top immediately on the receipt of this question paper. Use a separate answer book for each section. The figures in the right-hand margin indicate marks.}

\bigskip

\begin{center}
   \textbf{Section A : Environmental Biology} \\
   \textbf{(Marks : 45)}
\end{center}

\smallskip



\noindent   \textit{Note: Answer four questions including Question 1 which is compulsory.}

\bigskip



\noindent   \textbf{Question 1.} Answer the following parts:

\begin{parts}
  \item State whether the following statements are true or false. Give brief justification:\pts{$2  \times 5 = 10$}
  
\begin{parts}
    \item Rain water is the major source of fresh water resource.
    \item Standing crop amounts living material in different tropic level.
    \item Energy flow is a multidirectional process.
    \item Predation is a kind of interaction.
    \item Frequency, abundance and density are the qualitative characteristics of a community.
  \end{parts}

  \item Define the following:\pts{$1  \times 5 = 5$}
  
\begin{parts}
    \item Aquifers
    \item Keystone species
    \item Carrying capacity
    \item Land use planning
    \item Ecological density.
  \end{parts}
\end{parts}

\medskip


\noindent   \textbf{Question 2.} Describe the principle and methods of soil conservation.\pts{10}

\medskip


\noindent   \textbf{Question 3.} Discuss different level of biodiversity and also write about the main threats to biodiversity.\pts{10}

\medskip


\noindent   \textbf{Question 4.} Define air pollution. Describe its main sources, type of pollutants and its impact on the air pollution.\pts{10}

\medskip


\noindent   \textbf{Question 5.} Write an essay on population characteristics.\pts{10}

\pagebreak


\begin{center}
   \textbf{Section B : Systematics} \\
   \textbf{(Marks : 25)}
\end{center}

\smallskip



\noindent   \textit{Note: Answer three questions including Question 1 which is compulsory.}

\bigskip



\noindent   \textbf{Question 1.} Answer the following parts:

\begin{parts}
  \item Define the following:\pts{$1  \times 3 = 3$}
  
\begin{parts}
    \item Tautonym
    \item Gamma taxonomy
    \item ICZN.
  \end{parts}

  \item State whether the following statements are TRUE or FALSE giving reasons in 2--3 lines:\pts{$1  \times 4 = 4$}
  
\begin{parts}
    \item A paratype is a single specimen selected by the author of a species as its type.
    \item Tribe comes between order and class of the taxonomic category.
    \item Semicolon is always used in between author's name and the year.
    \item Typification is the grouping of objects or information based on similarities.
  \end{parts}
\end{parts}

\medskip


\noindent   \textbf{Question 2.} What do you understand by Zoological classification? Describe the various kinds of Zoological classification.\pts{9}

\medskip


\noindent   \textbf{Question 3.} Discuss various molecular techniques which are employed to identify different species of animals.\pts{9}

\medskip


\noindent   \textbf{Question 4.} Write notes on the following:\pts{$4.5  \times 2 = 9$}

\begin{parts}
  \item Morphological characters
  \item Biological species concept.
\end{parts}

\vfill
\rule{\linewidth}{0.4pt}

\begin{center}\small
\href{https://bhu-exam.vercel.app/}{Click Here}\\[4pt]\href{https://me-aryan.vercel.app/}{Click Here}\\[4pt]\href{https://quashan-me.vercel.app/}{Click Here}
\end{center}

\end{document}
